\documentclass[10pt,a4paper]{amsart}
\usepackage{amsmath,amssymb,amsthm,graphicx}
% Retained minimal preamble for the Stacks Project chapter flat.tex (More on Flatness).
% Verbatim excerpt of the chapter preamble: Xy-pic setup (native preamble lines 14-18) and the native
% theorem declarations with the mathematical macros (native preamble lines 154-217). Copied mechanically;
% nothing added, removed or redefined. The wrapper class and the amsmath/amssymb/amsthm/graphicx base
% packages are supplied by the assembler, so the class/externaldocument part of the native preamble is omitted.
% For commutative diagrams we use Xy-pic
\usepackage[all]{xy}

% We use 2cell for 2-commutative diagrams.
\xyoption{2cell}
% Theorem environments.
%
\theoremstyle{plain}
\newtheorem{theorem}[subsection]{Theorem}
\newtheorem{proposition}[subsection]{Proposition}
\newtheorem{lemma}[subsection]{Lemma}

\theoremstyle{definition}
\newtheorem{definition}[subsection]{Definition}
\newtheorem{example}[subsection]{Example}
\newtheorem{exercise}[subsection]{Exercise}
\newtheorem{situation}[subsection]{Situation}

\theoremstyle{remark}
\newtheorem{remark}[subsection]{Remark}
\newtheorem{remarks}[subsection]{Remarks}

\numberwithin{equation}{subsection}

% Macros
%
\def\lim{\mathop{\mathrm{lim}}\nolimits}
\def\colim{\mathop{\mathrm{colim}}\nolimits}
\def\Spec{\mathop{\mathrm{Spec}}}
\def\Hom{\mathop{\mathrm{Hom}}\nolimits}
\def\Ext{\mathop{\mathrm{Ext}}\nolimits}
\def\SheafHom{\mathop{\mathcal{H}\!\mathit{om}}\nolimits}
\def\SheafExt{\mathop{\mathcal{E}\!\mathit{xt}}\nolimits}
\def\Sch{\mathit{Sch}}
\def\Mor{\mathop{\mathrm{Mor}}\nolimits}
\def\Ob{\mathop{\mathrm{Ob}}\nolimits}
\def\Sh{\mathop{\mathit{Sh}}\nolimits}
\def\NL{\mathop{N\!L}\nolimits}
\def\CH{\mathop{\mathrm{CH}}\nolimits}
\def\proetale{{pro\text{-}\acute{e}tale}}
\def\etale{{\acute{e}tale}}
\def\QCoh{\mathit{QCoh}}
\def\Ker{\mathop{\mathrm{Ker}}}
\def\Im{\mathop{\mathrm{Im}}}
\def\Coker{\mathop{\mathrm{Coker}}}
\def\Coim{\mathop{\mathrm{Coim}}}

% Boxtimes
%
\DeclareMathSymbol{\boxtimes}{\mathbin}{AMSa}{"02}

%
% Macros for moduli stacks/spaces
%
\def\QCohstack{\mathcal{QC}\!\mathit{oh}}
\def\Cohstack{\mathcal{C}\!\mathit{oh}}
\def\Spacesstack{\mathcal{S}\!\mathit{paces}}
\def\Quotfunctor{\mathrm{Quot}}
\def\Hilbfunctor{\mathrm{Hilb}}
\def\Curvesstack{\mathcal{C}\!\mathit{urves}}
\def\Polarizedstack{\mathcal{P}\!\mathit{olarized}}
\def\Complexesstack{\mathcal{C}\!\mathit{omplexes}}
% \Pic is the operator that assigns to X its picard group, usage \Pic(X)
% \Picardstack_{X/B} denotes the Picard stack of X over B
% \Picardfunctor_{X/B} denotes the Picard functor of X over B
\def\Pic{\mathop{\mathrm{Pic}}\nolimits}
\def\Picardstack{\mathcal{P}\!\mathit{ic}}
\def\Picardfunctor{\mathrm{Pic}}
\def\Deformationcategory{\mathcal{D}\!\mathit{ef}}

% --- appended 2026-10-12 (Chief): the project own definitions of the reference/slogan/history
% display environments, copied mechanically from the shared Stacks preamble lines 8-12
% (\usepackage{verbatim} plus the three \newenvironment lines, which the chapter preamble keeps
% verbatim from the project preamble). Nothing added, removed or redefined: these environments
% carry no mathematical content and the project itself typesets them as comments. Without these
% three lines a counted statement whose own text uses them fails with Environment slogan undefined.
\usepackage{verbatim}
\newenvironment{reference}{\comment}{\endcomment}
\newenvironment{slogan}{\comment}{\endcomment}
\newenvironment{history}{\comment}{\endcomment}
\usepackage{cleveref}
\crefname{originallemma}{original lemma}{original lemmas}
\Crefname{originallemma}{Original lemma}{Original lemmas}
\makeatletter
\newlabel{lemma-finite-type-flat-at-point-free}{{05KY}{source}{}{sourcepointer-1}{}}
\newlabel{lemma-finite-type-flat-at-point-free@cref}{{[originallemma][0][]05KY}{[1][0][]source}{}{}{}}
\makeatother
\title{The pointwise freeness statement for a finite type module without the finite presentation hypothesis}
\author{The Stacks Project Authors}
\date{stacks-project:flat.tex:local-snapshot-20261011; GNU FDL 1.2 or later}
\begin{document}
\maketitle
\noindent\textit{Editorial note.} All delivered mathematical text is the author's own, copied from the pinned Stacks Project chapter file flat.tex; no statement or proof step was rewritten, reordered or removed.\par\medskip
\noindent\textit{Editorial note.} Preamble provenance: the chapter's own retained minimal preamble (sources/stacks-flat/retained-preamble.tex, lines 1-85) is a verbatim excerpt of the native Xy-pic setup and of the native theorem declarations and mathematical macros, followed by the project's own reference/slogan/history display-environment definitions copied mechanically from the shared preamble; the wrapper class and the amsmath/amssymb/amsthm/graphicx base packages are supplied by the assembler, and extra\_packages is left empty because everything needed is already in the retained preamble. No macro was invented or redefined.\par\medskip
\noindent\textit{Editorial note.} Same-chapter pointers are located at their \textbackslash{}label lines inside the pinned chapter file; cross-chapter pointers are located in the project's own tag list staged in the snapshot (sources/stacks-chow/tags.txt), at the verbatim lines carrying their canonical tags.\par\medskip
\noindent\textit{Allowed standard prerequisites (editorial).}\begin{itemize}
\item The pointwise freeness statement in the finite presentation case. Same chapter flat.tex, label lemma-finite-type-flat-at-point-free, invoked at proof line 3537. Applied to the ambient affine space over the base, the image point and the pushforward of the module, it produces the elementary etale neighbourhood, the affine etale neighbourhood and the freeness of the global sections, from which the solution for the original morphism is cut out.
\item The module pushed forward along a closed immersion has the same stalk over the base. Same chapter flat.tex, the remark at proof line 3534, using the identification of the stalks. Shows that the pushforward of the module to the affine space is still flat over the base at the image point, which is the hypothesis needed to apply the finite presentation statement there.
\end{itemize}\medskip
\ifdefined\hypertarget\hypertarget{sourcepointer-1}{}\fi
\noindent\textit{Source reference.} Original lemma 05KY: sources/stacks-flat/pinned-original-source.tex, source lines 3461--3461. Navigation pointer for \textbackslash{}ref\{lemma-finite-type-flat-at-point-free\} used inside the counted span (line 3527); the argument that consumes it is retained in the span.\par\medskip
\begin{lemma}
\label{lemma-finite-type-flat-at-point-free-variant}
Let $f : X \to S$ be a morphism of schemes.
Let $\mathcal{F}$ be a quasi-coherent sheaf on $X$.
Let $x \in X$ with image $s \in S$.
Assume that
\begin{enumerate}
\item $f$ is locally of finite type,
\item $\mathcal{F}$ is of finite type, and
\item $\mathcal{F}$ is flat at $x$ over $S$.
\end{enumerate}
Then there exists an elementary \'etale neighbourhood $(S', s') \to (S, s)$
and a commutative diagram of pointed schemes
$$
\xymatrix{
(X, x) \ar[d] & (X', x') \ar[l]^g \ar[d] \\
(S, s) & (\Spec(\mathcal{O}_{S', s'}), s') \ar[l]
}
$$
such that $X' \to X \times_S \Spec(\mathcal{O}_{S', s'})$
is \'etale, $\kappa(x) = \kappa(x')$, the scheme $X'$ is
affine, and such that $\Gamma(X', g^*\mathcal{F})$ is a free
$\mathcal{O}_{S', s'}$-module.
\end{lemma}

\begin{proof}
(The only difference with
Lemma \ref{lemma-finite-type-flat-at-point-free}
is that we do not assume $f$ is of finite presentation.)
The problem is local on $X$ and $S$. Hence we may assume $X$ and
$S$ are affine, say $X = \Spec(B)$ and $S = \Spec(A)$.
Since $B$ is a finite type $A$-algebra we can find a surjection
$A[x_1, \ldots, x_n] \to B$. In other words, we can choose a closed
immersion $i : X \to \mathbf{A}^n_S$. Set $t = i(x)$ and
$\mathcal{G} = i_*\mathcal{F}$. Note that $\mathcal{G}_t \cong \mathcal{F}_x$
are $\mathcal{O}_{S, s}$-modules. Hence $\mathcal{G}$ is flat over $S$ at $t$.
We apply
Lemma \ref{lemma-finite-type-flat-at-point-free}
to the morphism $\mathbf{A}^n_S \to S$, the point $t$, and the
sheaf $\mathcal{G}$. Thus we can find an
elementary \'etale neighbourhood $(S', s') \to (S, s)$
and a commutative diagram of pointed schemes
$$
\xymatrix{
(\mathbf{A}^n_S, t) \ar[d] & (Y, y) \ar[l]^h \ar[d] \\
(S, s) & (\Spec(\mathcal{O}_{S', s'}), s') \ar[l]
}
$$
such that $Y \to \mathbf{A}^n_{\mathcal{O}_{S', s'}}$
is \'etale, $\kappa(t) = \kappa(y)$, the scheme $Y$ is
affine, and such that $\Gamma(Y, h^*\mathcal{G})$ is a projective
$\mathcal{O}_{S', s'}$-module. Then a solution to the original
problem is given by the closed subscheme
$X' = Y \times_{\mathbf{A}^n_S} X$ of $Y$.
\end{proof}

\end{document}
